%=============================================================================
% WAVE 4, ITERATION 10: PATENT 1 + PATENT 8 --- DEAD PATENTS MONITOR
% Plus: Cross-Patent Error Checking and T4 Cold Spot Contribution
%
% Agent: P1+P8 Dead Patents Monitor (W4i10)
% Model: Opus 4.6
% Date: 20 March 2026
%
% MANDATE:
%   1. Confirm P1 and P8 still dead via 2025-2026 literature search
%   2. Cross-patent error checking of numerical claims in corpus
%   3. T4 Cold Spot contribution
%   4. Report TOP 5 findings ranked by impact
%
% INHERITED FACTS (not re-derived):
%   - P1 DEAD (W3i7 FINAL): signal ~0.81 zs, gap 10^15
%   - P8 DEAD (W3i7 FINAL): 1.6 Earth masses minimum at 1 AU
%   - Two-component: psi = psi_grav + delta_psi_EM
%   - M_eff = 3.01e6 kg (not 10^-6)
%   - Volume Cancellation Theorem, Passive Superiority Theorem
%   - SQMS bound: |lambda-1| < 1.56e-9
%=============================================================================

\documentclass[11pt]{article}
\usepackage[margin=2.5cm]{geometry}
\usepackage{amsmath,amssymb}
\usepackage{booktabs}
\usepackage{enumitem}
\usepackage{float}
\usepackage{xcolor}
\usepackage[most]{tcolorbox}
\usepackage{hyperref}

\newcommand{\ee}[1]{\times 10^{#1}}
\newcommand{\lamdiff}{|\lambda{-}1|}

\title{\textbf{Wave 4 Iteration 10: P1 + P8 Dead Patents Monitor}\\[0.5em]
\large 2025--2026 Literature Confirmation, Cross-Patent Error Audit,\\
and T4 Cold Spot Status Update}
\author{DFD Research Programme --- P1+P8 Agent, W4i10 (Opus 4.6)}
\date{20 March 2026}

\begin{document}
\maketitle

%=============================================================================
\section*{Executive Summary: TOP 5 FINDINGS (Ranked by Impact)}
%=============================================================================

\begin{tcolorbox}[colback=red!5!white, colframe=red!60!black,
  title=\textbf{TOP 5 FINDINGS --- W4i10}]

\textbf{Finding 1 (CRITICAL --- Corpus Inconsistency):} The
MASTER\_FINDINGS\_CORPUS states T4 Cold Spot is ``Ameliorated to 5--8x''
with ``tentatively $\sim$1--2x via void-evolution cancellation $\sim$73\%
(W3i7).'' However, W3i9 \emph{revoked} the W3i7 cancellation, finding
the 73\% cancellation was based on an incorrect treatment of $\dot{x}$
in the deep-MOND regime. W3i9 upgraded T4 from ``tentatively resolved''
to \textbf{``serious tension'' at 3--9x}. The corpus was \textbf{not
updated} to reflect this W3i9 correction. \textbf{The corpus T4 entry
is stale and misleading.}

\textbf{Finding 2 (SIGNIFICANT --- P8 Executive Summary Error):} The
W3i7 P6/P8/P9 executive summary states ``Minimum mass for 1~kbit/s at
1~AU: $2.7\ee{17}$~kg (Ceres-mass).'' The detailed derivation
(Eq.~(56) in W3i7) shows this is \textbf{wrong}: the 1~AU mass is
$9.4\ee{24}$~kg (1.6~$M_\oplus$). The $2.7\ee{17}$~kg figure is for
the submarine case at 500~m. The W3i8 convergence audit then
propagated this error by listing a ``correction'' from 1.6~$M_\oplus$
to $2.7\ee{17}$~kg as though one replaced the other. Both figures are
correct for their respective scenarios. \textbf{The executive summary
mislabels the submarine mass as the interplanetary mass, a factor of
$\sim 10^7$ error.}

\textbf{Finding 3 (MODERATE --- Wide Binary Prediction Under Pressure):}
DFD Prediction 10 claims $+42\%$ velocity enhancement at 10,000~AU.
Recent Gaia DR3 analyses (Banik et al.\ 2024; Pittordis \& Sutherland
2025) find Newtonian dynamics preferred at $19\sigma$ and MOND excluded
at $16\sigma$ in wide binaries. While some conflicting analyses exist
(Chae 2024), the most rigorous studies show \textbf{no evidence for
MOND-like enhancement}. This is not yet falsifying (systematic
uncertainties from triple contamination remain), but it constitutes
growing observational pressure on a core DFD prediction.

\textbf{Finding 4 (MODERATE --- DESI DR2 Updates):} DESI DR2
(published March 2025) strengthens evidence for dynamical dark energy
($w_0 > -1$, $w_a < 0$) with improved precision ($\sim 2\times$ over
DR1). DFD predicts opposite signs ($w_0 = -1.183$, $w_a = +0.183$).
The tension level should be reassessed with DR2 data, but the
sign mismatch persists. However, DESI's evidence for
$w \neq -1$ is qualitatively in DFD's direction (non-constant dark
energy), even if the specific parameterisation disagrees.

\textbf{Finding 5 (CONFIRMED --- P1 and P8 Remain Dead):} 2025--2026
literature search confirms no new physics changes the P1 or P8
assessments. Fifth force bounds have tightened (GRAVITY Galactic Centre:
$|\alpha| < 0.003$; neutron star cooling: strongest scalar-nucleon
limits). No laboratory MOND detection. SQMS/Dark SRF at Fermilab
continues to improve dark photon bounds (PRL March 2026, order of
magnitude improvement), but these constrain vector not scalar couplings.
P1 and P8 death assessments are reconfirmed.
\end{tcolorbox}


%=============================================================================
\part{P1 and P8 Death Confirmation}
%=============================================================================

\section{P1 (Field Drag): Still Dead}

\subsection{2025--2026 Literature Search Results}

Searched for: gravitational scalar fields, fifth force experiments,
screened scalar-tensor theories, MOND laboratory tests.

\begin{enumerate}[nosep]
\item \textbf{Atomic clock tests of screened scalar-tensor gravity}
  (Phys.\ Rev.\ D, March 2025): Proposes gravitational redshift
  measurements to test screened scalar-tensor models. Does not achieve
  new bounds relevant to the DFD coupling regime ($\lamdiff < 1.56\ee{-9}$).

\item \textbf{GRAVITY instrument at Galactic Centre} (2025): Upper limit
  $|\alpha| < 0.003$ for Yukawa fifth force at $\sim 200$~AU scale.
  Improvement of $\sim 10\times$ over previous. Constrains modified
  gravity at astrophysical scales but does not open new parameter space
  for P1 drag reduction, which requires $\lamdiff > 1.85\ee{-2}$.

\item \textbf{Neutron star cooling constraints} (December 2025): Most
  stringent limits on scalar-nucleon coupling from cooling rate analysis.
  Further tightens scalar field bounds. No relief for P1.

\item \textbf{OSIRIS-REx / asteroid Bennu constraints} (2024--2025):
  Constraints on fifth forces and ultralight dark matter at asteroid
  scales. No new parameter space for P1.
\end{enumerate}

\textbf{Verdict}: No new physics development in 2025--2026 changes
the P1 assessment. The fundamental barrier remains the $c^2/2$ lever
arm: generating detectable psi-gradients requires EM energy densities
exceeding $5.8\ee{16}$~J/m$^3$ ($10^{12}\times$ above SRF capability).
Signal remains $\sim 0.81$~zs, gap $10^{15}$.
\textbf{P1 is confirmed DEAD.}


\section{P8 (Communications): Still Dead}

\subsection{2025--2026 Literature Search Results}

Searched for: gravitational communication, scalar wave experiments,
psi-field transmission.

\begin{enumerate}[nosep]
\item \textbf{Gravitational Communication survey} (arXiv 2501.03251,
  January 2025): Comprehensive review of gravitational wave
  communication concepts. Confirms that all proposed mechanisms require
  planetary-scale masses or beyond. No breakthrough modulation technique
  identified.

\item \textbf{High-frequency GW detection} (October 2025): Table-top
  interferometric system achieves $10^{-19}$~m sensitivity at high
  frequencies. While impressive, this is a \emph{detection} advance,
  not a \emph{generation} advance. The P8 bottleneck is transmitter
  mass, not receiver sensitivity.

\item \textbf{Dark SRF at Fermilab} (PRL, March 2026): Order-of-magnitude
  improvement in dark photon exclusion from refined SRF analysis. Relevant
  to P11/P2 detection programme, but does not change P8 mass requirements.
\end{enumerate}

\textbf{Verdict}: No new physics development changes the P8 assessment.
The minimum transmitter mass for 1~kbit/s remains:
\begin{itemize}[nosep]
\item At 1~AU (interplanetary): $9.4\ee{24}$~kg $= 1.6\,M_\oplus$
\item At 500~m (submarine): $3\ee{17}$~kg $\approx$ Ceres-mass
\end{itemize}
Both are technologically inaccessible by $\geq 10^{11}\times$.
$Q_\psi$ amplification does not help (W3i7 Theorem).
\textbf{P8 is confirmed DEAD.}


%=============================================================================
\part{Cross-Patent Numerical Error Audit}
%=============================================================================

\section{Finding 1: Corpus T4 Entry Is Stale (CRITICAL)}
\label{sec:T4-stale}

The MASTER\_FINDINGS\_CORPUS (compiled 20 March 2026, labelled ``FINAL
corpus incorporating Wave 3 iterations 1--8'') contains the following
T4-related statements:

\begin{quote}
``Cold Spot ISW was initially 21x overpredicted (W3i3). Successive
iterations reduced this: 10--15x (W3i5), 3--5x (W3i6), tentatively
$\sim$1--2x via void-evolution cancellation $\sim$73\% (W3i7). This
remains the only genuine quantitative tension at 5--8x (conservative
estimate).''
\end{quote}

And in the tensions table:
\begin{quote}
``T4 (Cold Spot) | Ameliorated to 5--8x | \textbf{MODERATE}''
\end{quote}

However, the W3i9 document (W3i9\_T4\_ColdSpot.tex) performed a
definitive reanalysis and found:

\begin{enumerate}[nosep]
\item The W3i7 void-evolution cancellation of 73\% is \textbf{invalid}.
  It used an incomplete $\dot{x}$ expression that omitted the background
  dilution term ($-3H/2$) in the deep-MOND regime where $x \propto
  \sqrt{g_N}$.
\item With the corrected treatment, void evolution either has no net
  effect ($f_{\rm MOND} = 3$) or \emph{amplifies} the ISW signal
  ($f_{\rm MOND} < 3$).
\item T4 status was upgraded from ``tentatively resolved'' to
  \textbf{``serious tension''} at 3--9$\times$ overprediction.
\item The most promising resolution pathway is environmental MOND
  screening via Hubble-flow acceleration at void scales.
\end{enumerate}

\textbf{Impact}: The corpus is internally inconsistent --- it claims
``iterations 1--8'' but the T4 entry contradicts the W3i9 finding. The
corpus must be updated to reflect the W3i9 result:

\begin{center}
\begin{tabular}{lcc}
\toprule
\textbf{Item} & \textbf{Current Corpus} & \textbf{Should Be} \\
\midrule
T4 status & Ameliorated 5--8x & \textbf{Serious tension 3--9x} \\
Void-evolution & 73\% cancellation & \textbf{Cancellation invalid} \\
T4 severity & Moderate & \textbf{Serious} \\
\bottomrule
\end{tabular}
\end{center}


\section{Finding 2: P8 Executive Summary Mass Error (SIGNIFICANT)}
\label{sec:P8-mass}

The W3i7 P6/P8/P9 executive summary states:

\begin{quote}
``Minimum mass for 1~kbit/s at 1~AU: $2.7\ee{17}$~kg (Ceres-mass).
Psi-comms remains technologically infeasible by $\sim 10^{11}\times$.''
\end{quote}

The detailed derivation in the same document (Section~6.2,
Eq.~(56)) shows:
\begin{itemize}[nosep]
\item At 1~AU with $N = 10^6$ sensors: $M_s = 9.4\ee{24}$~kg $= 1.6\,M_\oplus$
\item At 500~m (submarine) with $N = 10^4$ sensors: $M_s = 3\ee{17}$~kg $\approx$ Ceres
\end{itemize}

The executive summary conflates the two, labelling the 500~m submarine
figure as the 1~AU interplanetary figure. This is a factor of
$\sim 3.5\ee{7}$ error in the executive summary.

The W3i8 convergence audit then listed ``P8 minimum transmitter mass:
$1.6\,M_\oplus \to 2.7\ee{17}$~kg (Ceres-mass)'' as though one
corrected the other. In fact, both are valid for different scenarios.

\textbf{Consequence}: The corpus correctly states 1.6 Earth masses for
1~AU. The bottom-line P8 death verdict is unaffected (both masses are
inaccessible). But the W3i7 executive summary and W3i8 convergence
audit contain a misleading cross-reference.

\textbf{Recommended correction}: The W3i7 executive summary should read:
``Minimum mass for 1~kbit/s at 1~AU: $\sim 10^{25}$~kg
(1.6~$M_\oplus$). For submarine at 500~m: $\sim 3\ee{17}$~kg
(Ceres-mass). Mass gap $\geq 10^{11}\times$ in all cases.''


\section{Finding 3: Wide Binary Prediction Under Observational Pressure}
\label{sec:wide-binary}

DFD Prediction 10 (Section 3, Tier 1 of the corpus):
\begin{quote}
``Wide Binary Enhancement +42\% at 10,000 AU ($M_\odot$) using
corrected simple-mu formula.''
\end{quote}

Recent Gaia DR3 wide binary analyses:
\begin{itemize}[nosep]
\item \textbf{Banik et al.\ (2024, MNRAS 527)}: Comprehensive
  likelihood analysis finds Newtonian gravity preferred at $19\sigma$
  over standard MOND. MOND excluded at $16\sigma$.
\item \textbf{Pittordis \& Sutherland (2025, MNRAS 547)}: Updated
  quality framework with stringent contamination cuts. No evidence for
  MOND-induced velocity boosts. Meta-analysis shows apparent MOND signals
  diminish with improved methodology.
\item \textbf{Chae (2024)}: Claims evidence for MOND-like signal with
  different sample selection. Contradicts Banik/Pittordis.
\end{itemize}

The DFD prediction uses $\mu(x) = x/(1+x)$, which gives a $+42\%$
enhancement at separations where $g \sim a_0$. The Banik/Pittordis
results, if robust, rule out this enhancement at $>10\sigma$.

\textbf{Assessment}: This is not yet a clear-cut falsification because:
(a) the Chae analysis disagrees; (b) triple-star contamination remains
a dominant systematic; (c) the external field effect in DFD modifies
the prediction at the Solar neighbourhood level. However, the trend of
increasingly rigorous analyses finding no MOND signal is concerning.

\textbf{Recommendation}: The corpus should flag Prediction 10 with a
caution note referencing the 2024--2025 Gaia DR3 results.


\section{Finding 4: DESI DR2 Tension Update}
\label{sec:DESI}

DESI DR2 (March 2025) provides a factor of $\sim 2$ improvement in BAO
precision. Key findings relevant to DFD:
\begin{itemize}[nosep]
\item Evidence for dynamical dark energy ($w \neq -1$) has strengthened.
\item DESI+CMB+SNe prefer $w_0 > -1$ and $w_a < 0$.
\item DFD predicts $w_0 = -1.183$ and $w_a = +0.183$ (opposite signs on
  both parameters).
\end{itemize}

The DFD--DESI sign mismatch persists. The tension level should be
re-evaluated with DR2 data, but the qualitative picture is unchanged
from the corpus's 2.4$\sigma$ estimate. The tension may have increased
slightly with improved DESI precision.

\textbf{Note}: DESI's evidence for $w \neq -1$ is qualitatively
favourable to DFD (non-constant dark energy), even though the specific
CPL parameterisation disagrees.


\section{Numerical Spot-Checks (No Errors Found)}
\label{sec:spot-checks}

The following key numerical claims were independently verified:

\begin{center}
\begin{tabular}{lcc}
\toprule
\textbf{Claim} & \textbf{Corpus Value} & \textbf{Verified} \\
\midrule
$a_0 = 2\sqrt{\alpha}\,cH_0$ & $1.12\ee{-10}$~m/s$^2$ & $\checkmark$ \\
Shadow ratio $2e/(3\sqrt{3})$ & 1.046 ($+4.6\%$) & $\checkmark$ \\
SQMS bound & $\lamdiff < 1.56\ee{-9}$ & Consistent across all refs \\
$\kappa = 3/(8\alpha)$ & $\approx 51.4$ & $\checkmark$ \\
Formalism $a_0$ rounding & $1.2\ee{-10}$ vs $1.12\ee{-10}$ & Minor
  rounding, not error \\
\bottomrule
\end{tabular}
\end{center}


%=============================================================================
\part{T4 Cold Spot: Status Contribution}
%=============================================================================

\section{Current State After W3i9}

The W3i9 definitive analysis established:
\begin{enumerate}[nosep]
\item Static DFD ISW: $-350$ to $-700\,\mu$K (depending on
  $\delta_v = -0.10$ to $-0.20$).
\item Void-evolution: amplifies (does not cancel) the signal for
  $f_{\rm MOND} < 3$.
\item Best-case overprediction: $\sim 2\times$ ($\delta_v = -0.10$,
  0\% primordial, no amplification).
\item Canonical overprediction: $\sim 5\times$ ($\delta_v = -0.14$,
  50\% primordial).
\item Most promising resolution: environmental MOND screening via
  Hubble-flow acceleration.
\end{enumerate}

\section{2025--2026 Cold Spot Observations}

Literature search for ``Cold Spot CMB ISW Eridanus supervoid 2025--2026'':

\begin{enumerate}[nosep]
\item The DES Year-3 confirmation (Kov\'acs et al.\ 2022, MNRAS 510)
  remains the most recent detailed characterisation. The Eridanus
  supervoid is confirmed as a significant underdensity at $z < 0.2$,
  1.8~Gly wide, $\sim 30\%$ underdense.
\item Standard $\Lambda$CDM ISW from the supervoid accounts for only
  10--20\% of the observed temperature decrement.
\item CMB lensing amplitude at the Cold Spot centre is $\sim 30\%$
  lower than expected from $\Lambda$CDM simulations.
\item No new 2025--2026 spectroscopic mapping has been published.
  DESI DR2 spectroscopic characterisation is expected 2027--2028.
\end{enumerate}

\textbf{No new data changes the T4 assessment.} The decisive observation
remains DESI spectroscopic mapping of the void depth.

\section{Potential Resolution Pathway: Hubble-Flow Screening}

The W3i9 suggestion that Hubble-flow acceleration at void scales
($H_0^2 R \sim 0.5\,a_0$) could provide environmental screening
deserves theoretical development. At void scale $R \sim 300$~Mpc:
\begin{equation}
\frac{H_0^2\,R}{a_0} = \frac{(2.2\ee{-18})^2 \times 9.3\ee{24}}{1.12\ee{-10}}
\approx 0.4.
\end{equation}
This would place the void in the $\mu$-transition regime
($x_{\rm eff} \sim 0.4$, $\mu \sim 0.29$) rather than deep-MOND
($x \sim 10^{-3}$, $\mu \sim 10^{-3}$). The ISW enhancement would
drop from $\sim 100\times$ to $\sim 3\times$, potentially resolving T4.

However, this requires modifying the MOND interpolation to include
cosmological accelerations, which:
\begin{itemize}[nosep]
\item Has theoretical precedent (Milgrom's ``external field effect'')
\item But at cosmological scales introduces ambiguity about what
  constitutes the ``external'' field
\item Would need to be shown consistent with galaxy-scale MOND
  phenomenology (where it must not apply)
\end{itemize}

\textbf{This remains the single most important theoretical problem
in DFD cosmology.}


%=============================================================================
\part{Summary and Recommendations}
%=============================================================================

\begin{tcolorbox}[colback=green!5!white, colframe=green!60!black,
  title=\textbf{W4i10 Summary}]

\textbf{P1}: Confirmed DEAD. No 2025--2026 developments change the
assessment. Fifth force bounds continue to tighten.

\textbf{P8}: Confirmed DEAD. No breakthrough modulation technique.
Executive summary mass error identified (Finding 2) but does not
affect the death verdict.

\textbf{Corpus inconsistencies found}: 2 significant.
\begin{enumerate}[nosep]
\item T4 entry not updated from W3i9 (serious tension, not
  ameliorated). \textbf{CRITICAL.}
\item P8 mass figure cross-reference error in W3i7/W3i8.
  \textbf{SIGNIFICANT but non-impacting on verdict.}
\end{enumerate}

\textbf{New observational pressure}: Wide binary Gaia DR3 analyses
(Banik 2024, Pittordis 2025) challenge Prediction 10 ($+42\%$
enhancement). Not yet falsifying but concerning.

\textbf{DESI DR2}: Sign mismatch persists. Tension likely unchanged
or slightly increased with improved data.

\textbf{T4}: Remains the single most serious challenge. Environmental
MOND screening via Hubble-flow acceleration is the most promising
resolution pathway and needs theoretical development.
\end{tcolorbox}


\section*{Action Items for Corpus Maintainer}

\begin{enumerate}
\item \textbf{UPDATE} T4 entry in MASTER\_FINDINGS\_CORPUS: change
  ``Ameliorated to 5--8x'' to ``Serious tension 3--9x (W3i9)'';
  remove reference to 73\% void-evolution cancellation; update
  tensions table from ``Moderate'' to ``Serious.''
\item \textbf{FLAG} Prediction 10 (wide binaries) with caution note
  referencing Banik et al.\ (2024) and Pittordis \& Sutherland (2025).
\item \textbf{CLARIFY} P8 mass figures in Section 1 of corpus:
  distinguish 1~AU ($1.6\,M_\oplus$) from submarine ($\sim$Ceres)
  explicitly.
\item \textbf{REASSESS} DESI tension with DR2 data when full DFD
  Boltzmann code is available.
\end{enumerate}


\end{document}
